\section{A Four-Axis Taxonomy}
\label{sec:taxonomy}

\todo{This section defines the taxonomy axes. Each axis should be
illustrated with examples from the surveyed systems.}

We organize the landscape along four axes (\Cref{fig:taxonomy}).

% TODO: Create taxonomy figure
% \begin{figure}[t]
%     \centering
%     \includegraphics[width=\columnwidth]{figures/taxonomy-axes.pdf}
%     \caption{Four-axis taxonomy for AI agent memory systems.}
%     \label{fig:taxonomy}
% \end{figure}

\subsection{Axis 1: Memory Structure}

Following CoALA~\cite{sumers2024coala}, we classify memory modules as:

\begin{itemize}[leftmargin=*, nosep]
    \item \textbf{Working memory:} In-context information the LLM
    processes directly. Managed by the agent, not the memory system.
    \item \textbf{Episodic memory:} Raw experiences preserved with
    full context (who, when, where, about what).
    \item \textbf{Semantic memory:} Distilled, decontextualized
    knowledge extracted from episodes.
    \item \textbf{Procedural memory:} Executable skills and
    behavioral patterns.
    \item \textbf{Implicit memory:} Preferences and habits that
    emerge from accumulated interaction without explicit extraction.
    Not present in CoALA; we propose it as an extension.
\end{itemize}

\todo{Add table: distribution of memory types across all 90+ systems.
Show that most systems support only 1--2 types.}

\subsection{Axis 2: Retrieval Architecture}

Following Gao~\etal~\cite{gao2023rag}, we classify retrieval as:

\begin{itemize}[leftmargin=*, nosep]
    \item \textbf{Naive RAG:} Single retrieval signal (vector or
    keyword), direct injection. Examples: LangChain, Engram.
    \item \textbf{Advanced RAG:} Multiple signals or reranking.
    Examples: Generative Agents (recency $\times$ importance $\times$
    relevance), Mem0 (vector + graph + priority).
    \item \textbf{Modular RAG:} Composable pipeline with independent
    retrieval, fusion, and reranking stages. Examples: Alaya (BM25 +
    vector + graph + RRF), Zep/Graphiti (cosine + BM25 + graph + RRF
    + MMR).
\end{itemize}

\subsection{Axis 3: Memory Lifecycle}

Following Zhang~\etal~\cite{zhang2024survey} and
Hu~\etal~\cite{hu2025memory}:

\begin{itemize}[leftmargin=*, nosep]
    \item \textbf{Formation:} How memories enter the system (direct
    storage, LLM extraction, agent-directed).
    \item \textbf{Consolidation:} How episodic memories transform into
    semantic knowledge (CLS replay, LLM reflection, sleep-time
    processing).
    \item \textbf{Forgetting:} How memories weaken or disappear
    (no forgetting, single-curve decay, dual-strength decay, RL-learned).
    \item \textbf{Transformation:} How the memory store itself
    restructures (deduplication, contradiction resolution, pruning).
\end{itemize}

\subsection{Axis 4: Preference Emergence}

We propose this as a new axis not covered by existing surveys:

\begin{itemize}[leftmargin=*, nosep]
    \item \textbf{No preferences:} Most systems.
    \item \textbf{LLM-extracted profiles:} Mem0, Supermemory, Memobase.
    One-shot extraction, brittle.
    \item \textbf{Explicit opinion tracking:} Hindsight. Confidence-scored
    beliefs that evolve.
    \item \textbf{Emergent crystallization:} Alaya. Preferences crystallize
    from accumulated impressions without LLM involvement.
\end{itemize}
